\subsection{A Finite Alphabet Is Not a Peano System}

\begin{topicbox}{A Failed Carrier}
\begin{exposition}
Finite chains help isolate what the Peano axioms forbid. A finite alphabet can
support a partial next-letter operation, or it can support a cyclic successor
operation, but neither choice gives a Peano system.
\end{exposition}

\begin{example*}[The Finite Alphabet Fails]
Let
\[
P=\{A,B,C,\ldots,Z\}.
\]
If one tries to define
\[
S(A)=B,\quad S(B)=C,\quad \ldots,
\]
then there is no value of $S(Z)$ inside the same finite alphabet that extends
the one-way chain. If instead one defines $S(Z)=A$, then $S:P\to P$ is total
but the successor operation forms a cycle.
\end{example*}

\begin{remark*}[Failure modes]
The finite alphabet fails in one of two ways:
\[
\begin{aligned}
&\text{either } S \text{ is not a total map } P\to P,\\
&\text{or } S \text{ loops back into a cycle.}
\end{aligned}
\]
The first failure violates closure of successor. The second failure violates
the one-way successor structure forced by the Peano axioms.
\end{remark*}

\begin{remark*}[Interpretation]
A finite alphabet cannot be a Peano system. Peano systems are necessarily
infinite because induction and successor closure generate an unending chain
from the distinguished first element.
\end{remark*}
\end{topicbox}

